\documentclass{article}
\title{Practice Problem 2.46}
\author{CSSAPP}
\usepackage{booktabs, indentfirst, amsmath}

\begin{document}
\maketitle

\begin{flushleft}

    The imprecision of floating-point arithmetic can have disastrous effects.
    On February $25, 1991$, during the first Gulf War, an American Patriot Missile battery in
    Dharan, Saudi Arabia, failed to intercept an incoming Iraqi Scud missile. The Scud struck an
    American Army barracks and killed $28$ soldiers. The US General Accounting Office (GAO)
    conducted a detailed analysis of the failure[76] and determined that the underlying cause was
    an imprecision in a numeric calculation. In this exercise, you will reproduce part of the
    GAO's analysis.

    \noindent The Patriot system contains an internal clock, implemented as a counter that is
    incremented every $0.1$ seconds. To determine the time in seconds, the program will multiply
    the value of this counter by a $24$-bit quantity that was a fractional binary approximation to
    $\frac{1}{10}$. In particular, the binary representation of $\frac{1}{10}$ is the
    nonterminating sequence $0.000110011[0011]\ldots_2$, where the portion in the brackets is
    repeated indefinitely. The program approximated $0.1$, as a value $x$, by considering just the
    first $23$ bits of the sequence to the right of the binary point:
    \[
        x = 0.00011001100110011001100_2 \quad \text{(ref P.P.\ 2.51)}
    \]

    \smallbreak
    \noindent \textbf{A.} What is the binary representation of $0.1 - x$?\\
    \noindent \textbf{B.} What is the approximate decimal value of $0.1 - x$?\\
    \noindent \textbf{C.} The clock starts at $0$ when the system is first powered up and keeps
    counting up from there. In this case, the system had been running for about $100$ hours.
    What is the difference between the actual time and the time computed by the software?\\
    \noindent \textbf{D.} The system predicts where an incoming missile will appear based on its
    velocity and time of the last radar detection. Given that the Scud travels at around $2{,}000$
    meters per second, how far off was its prediction?

    \bigskip
    \hrule
    \bigskip

    \section*{Solution}

    \noindent \textbf{Setup.}
    The exact binary expansion of $\frac{1}{10}$ is
    \[
        0.1 = 0.0001\overline{1001}_2 = 0.000110011001100110011001100\ldots_2
    \]
    The Patriot system truncates this to $23$ bits after the binary point:
    \[
        x = 0.00011001100110011001100_2
    \]
    Note that $x$ agrees with the true value of $0.1$ exactly for the first $23$ bits;
    the entire error comes from the bits that were chopped off.

    \bigskip
    \noindent \textbf{Part A: Binary representation of $0.1 - x$}

    \smallskip
    \noindent Since the leading $23$ bits cancel, subtracting leaves only the
    remaining (``tail'') bits of $0.1$, shifted right by $23$ places:
    \[
        0.1 - x \;=\; \underbrace{0.000\ldots000}_{23 \text{ zeros}}1100110011\overline{0011}\ldots_2
    \]
    Equivalently,
    \[
        0.1 - x = 2^{-23} \times 0.\overline{1100}_2
    \]

    \bigskip
    \noindent \textbf{Part B: Decimal value of $0.1 - x$}

    \smallskip
    \noindent A repeating $4$-bit binary block $b$ has closed-form value
    $\dfrac{b}{2^4 - 1}$. Here $b = 1100_2 = 12$, so
    \[
        0.\overline{1100}_2 = \frac{12}{15} = 0.8
    \]
    Therefore
    \[
        0.1 - x = 0.8 \times 2^{-23} = 0.8 \times 1.1920929\times 10^{-7}
        \approx 9.54 \times 10^{-8}
    \]
    So each clock tick (each $0.1$-second interval) undercounts time by
    approximately \textbf{95 nanoseconds}.

    \bigskip
    \noindent \textbf{Part C: Accumulated clock error after 100 hours}

    \smallskip
    \noindent Number of $0.1$-second ticks in $100$ hours:
    \[
        100 \text{ hr} \times 3600 \text{ s/hr} = 360{,}000 \text{ s}
        \quad\Longrightarrow\quad
        \frac{360{,}000}{0.1} = 3{,}600{,}000 \text{ ticks}
    \]
    Total timing error:
    \[
        3{,}600{,}000 \times 9.54\times 10^{-8}\text{ s} \approx 0.343 \text{ s}
    \]
    The software clock lags the true time by approximately
    \textbf{0.343 seconds} after 100 hours of continuous operation.

    \bigskip
    \noindent \textbf{Part D: Resulting position error}

    \smallskip
    \noindent Using a missile speed of approximately $2000$ m/s:
    \[
        \text{Distance error} = v \times \Delta t \approx 2000 \text{ m/s} \times 0.343 \text{ s}
        \approx 686 \text{ m}
    \]
    The system's predicted intercept point was off by roughly
    \textbf{687 meters} --- enough that the incoming Scud fell outside the
    interceptor's engagement range.

    \bigskip
    \noindent \textbf{Summary}

    \begin{center}
        \begin{tabular}{@{}cl@{}}
            \toprule
            Part & Result \\
            \midrule
            A & $0.1-x = 0.000\ldots0\,1100110011\ldots_2$ (23 zeros, then repeating $1100$) \\
            B & $\approx 9.54\times10^{-8}$ \\
            C & $\approx 0.343$ s clock drift after 100 hrs \\
            D & $\approx 687$ m position error \\
            \bottomrule
        \end{tabular}
    \end{center}

\end{flushleft}
\end{document}